\documentclass[11pt]{article}

\usepackage[letterpaper,margin=1in]{geometry}
\usepackage{amsmath,amssymb}
\usepackage{booktabs}
\usepackage{longtable}
\usepackage{array}
\usepackage{xcolor}
\usepackage{listings}
\usepackage{url}
\usepackage{hyperref}
\usepackage{titlesec}
\usepackage{parskip}
\usepackage{fancyhdr}
\usepackage{enumitem}

\hypersetup{
  colorlinks=true,
  linkcolor=black,
  urlcolor=blue!50!black,
  citecolor=black,
  pdftitle={LonghornSilicon Reference Model API},
  pdfauthor={LonghornSilicon},
  pdfsubject={C++ / C / Python API reference for the ACU blocks}
}

\titleformat{\section}{\normalfont\Large\bfseries}{\thesection.}{0.6em}{}
\titleformat{\subsection}{\normalfont\large\bfseries}{\thesubsection}{0.6em}{}
\titleformat{\subsubsection}{\normalfont\normalsize\bfseries}{\thesubsubsection}{0.6em}{}

\pagestyle{fancy}
\fancyhf{}
\lhead{\small Reference Model API}
\rhead{\small Blocks 1 \& 2}
\cfoot{\thepage}
\renewcommand{\headrulewidth}{0.4pt}

\definecolor{codebg}{gray}{0.97}
\definecolor{codekey}{rgb}{0.13,0.29,0.53}
\definecolor{codecom}{rgb}{0.25,0.5,0.25}
\definecolor{codestr}{rgb}{0.5,0.13,0.13}
\lstset{
  basicstyle=\ttfamily\small,
  backgroundcolor=\color{codebg},
  frame=single,
  framerule=0.4pt,
  rulecolor=\color{black!30},
  xleftmargin=0.5em,
  xrightmargin=0.5em,
  showstringspaces=false,
  breaklines=true,
  keywordstyle=\color{codekey}\bfseries,
  commentstyle=\color{codecom}\itshape,
  stringstyle=\color{codestr},
  numbers=none
}

\setlist{itemsep=2pt,topsep=3pt}

\begin{document}

%==============================================================================
\thispagestyle{empty}
\begin{center}
\vspace*{1.2in}
{\Huge\bfseries Reference Model API}\\[0.4em]
{\LARGE\bfseries C++ Class \textperiodcentered\ extern "C" \textperiodcentered\ Python}\\[1.5em]

{\large \textsc{LonghornSilicon ACU --- Blocks 1 \& 2}}\\[2em]

\rule{0.7\textwidth}{0.4pt}\\[1em]

\begin{tabular}{rl}
\textbf{Date}:        & 2026-05-14 \\
\textbf{Covers}:      & Precision Controller, MAC Array \\
\textbf{Languages}:   & C++17, C99, Python 3.10+ \\
\textbf{Build}:       & \texttt{make} in \texttt{sw/reference\_model/} \\
\end{tabular}\\[1em]
\rule{0.7\textwidth}{0.4pt}

\vfill
\begin{minipage}{0.85\textwidth}
\small
\textbf{Purpose.} Formal API reference for the bit-accurate
reference implementations of the LonghornSilicon ACU blocks.
Each block ships three language bindings (C++ class, extern "C",
Python); all three pass identical test vectors. A compiler backend
chooses whichever binding fits its build system; numerical results
are guaranteed identical.
\end{minipage}

\vspace{0.3in}
\end{center}

\clearpage

\tableofcontents
\clearpage

%==============================================================================
\section{Quick Start}
\label{sec:quickstart}

\begin{lstlisting}[basicstyle=\ttfamily\small]
cd sw/reference_model
make test-all      # build C++ tests, run them, run all Python tests
make integration   # end-to-end attention-tile worked example
\end{lstlisting}

\noindent
Requires Python 3.10+, NumPy, and a C++17 compiler.

%==============================================================================
\section{Precision Controller}
\label{sec:pc}

\subsection{C++ Class}
\begin{lstlisting}[language=C++]
#include "precision_controller_ref.hpp"

lhsi::PrecisionController pc;
std::vector<int32_t> scores(/* 4096 INT8 values */);
bool decision = pc.process_tile(scores);   // true = FP16, false = INT8
\end{lstlisting}

\subsection{extern "C" API}
\begin{lstlisting}[language=C]
#include "precision_controller_ref.hpp"

lhsi_pc_handle_t* h = lhsi_pc_create();
int32_t scores[4096] = { /* ... */ };
int decision = lhsi_pc_process_tile(h, scores, 4096);
lhsi_pc_destroy(h);
\end{lstlisting}

Stateless variant (no handle):
\begin{lstlisting}[language=C]
int decision = lhsi_pc_decide(scores, 4096);
\end{lstlisting}

\subsection{Python}
\begin{lstlisting}[language=Python]
from precision_controller_ref import PrecisionController

pc = PrecisionController()
decision = pc.process_tile(scores)
# Or stateless:
decision = PrecisionController.decide(scores)
\end{lstlisting}

\subsection{Streaming API}
\label{sec:pc-stream}

For driver code or cycle-accurate simulators that want to issue one
score per clock (mirrors the chip's AXI-Stream interface).

\begin{lstlisting}[language=C++]
pc.reset();
for (size_t i = 0; i < scores.size(); ++i) {
    pc.tick(/*s_valid=*/true, scores[i], /*s_last=*/(i == scores.size() - 1));
    if (pc.d_valid()) {
        bool decision = pc.d_fp16();
    }
}
\end{lstlisting}

The Python class has the same shape:
\texttt{pc.tick(s\_valid=True, s\_data=x, s\_last=False)}.

%==============================================================================
\section{MAC Array}
\label{sec:mac}

\subsection{C++ Class}
\begin{lstlisting}[language=C++]
#include "mac_array_ref.hpp"

lhsi::mac::MacArray mac;

// INT8 path
std::vector<int8_t> A(M*K), B(K*N);
std::vector<int32_t> C(M*N);
mac.matmul_int8(A.data(), B.data(), C.data(), M, K, N);

// FP16 path (float storage; fp16 rounding on output)
std::vector<float> Af(M*K), Bf(K*N), Cf(M*N);
mac.matmul_fp16(Af.data(), Bf.data(), Cf.data(), M, K, N);

// Cost estimate for the compiler's scheduler
auto cost = mac.estimate(M, K, N, lhsi::mac::MacArray::DType::Int8);
// cost.cycles, cost.energy_pj
\end{lstlisting}

\subsection{extern "C" API}
\begin{lstlisting}[language=C]
#include "mac_array_ref.hpp"

int8_t  A[M*K], B[K*N];
int32_t C[M*N];
lhsi_mac_matmul_int8(A, B, C, M, K, N);

float Af[M*K], Bf[K*N], Cf[M*N];
lhsi_mac_matmul_fp16(Af, Bf, Cf, M, K, N);

lhsi_mac_cost_t cost;
lhsi_mac_estimate(M, K, N, LHSI_MAC_DTYPE_INT8, &cost);
\end{lstlisting}

\subsection{Python}
\begin{lstlisting}[language=Python]
import numpy as np
from mac_array_ref import MacArray

mac = MacArray()

# INT8
A = np.random.randint(-128, 128, (M, K), dtype=np.int8)
B = np.random.randint(-128, 128, (K, N), dtype=np.int8)
C = mac.matmul_int8(A, B)        # numpy int32 array

# FP16
A = np.random.uniform(-1, 1, (M, K)).astype(np.float16).astype(np.float32)
B = np.random.uniform(-1, 1, (K, N)).astype(np.float16).astype(np.float32)
C = mac.matmul_fp16(A, B)        # float32 storage, fp16-rounded output

# Cost
cost = mac.estimate(M, K, N, dtype="int8")
print(cost.cycles, cost.energy_pj)
\end{lstlisting}

\subsection{FP16 Utility Helpers}
\label{sec:mac-fp16-utils}

Exposed for boundaries where you need explicit fp16 rounding without
going through a matmul.

\begin{lstlisting}[language=C++]
float    rounded = lhsi::mac::round_to_fp16(x);
uint16_t bits    = lhsi::mac::float_to_fp16_bits(x);
float    back    = lhsi::mac::fp16_bits_to_float(bits);
\end{lstlisting}

Same functions are available in Python:

\begin{lstlisting}[language=Python]
from mac_array_ref import (
    round_to_fp16, float_to_fp16_bits, fp16_bits_to_float,
)
\end{lstlisting}

\clearpage

%==============================================================================
\section{Numerical Semantics (Frozen)}
\label{sec:semantics}

\subsection{Precision Controller}

Pure integer arithmetic in unsigned fixed-width modular form:

\begin{itemize}
\item Input \texttt{s\_data}: \texttt{SCORE\_WIDTH}-bit signed
      two's complement.
\item \texttt{max\_acc}: \texttt{SCORE\_WIDTH} bits (unsigned absolute).
\item \texttt{sum\_acc}: \texttt{SCORE\_WIDTH} $+ \log_2(N)$ bits
      (unsigned).
\item $\text{LHS} = \text{max} \ll \log_2(N)$ (free shift; wire
      routing).
\item $\text{RHS} = (\text{sum} \ll 3) + (\text{sum} \ll 1) =
      \text{sum} \times 10$.
\item Decision: $\text{LHS} > \text{RHS} \Rightarrow$ FP16.
\end{itemize}

If the reference disagrees with the chip on any input, the reference
is wrong. File an issue with the failing tile from
\texttt{rtl/tb/testvectors/}.

\subsection{MAC Array, INT8 Path}

Bit-exact integer matmul, \texttt{int32\_t} accumulator, no
saturation. The compiler is responsible for any requantization after
the matmul.

\subsection{MAC Array, FP16 Path}

\texttt{float} internally; IEEE-754 binary16 round-to-nearest-even
on each output element. See \path{docs/mac_array_design.pdf}~\S3.2
for the v0.1 simplifications and what may change in v0.2.

%==============================================================================
\section{Linking Options}
\label{sec:linking}

\subsection{Static Library}
\begin{lstlisting}[basicstyle=\ttfamily\small]
make static    # produces libprecision_controller_ref.a + libmac_array_ref.a
\end{lstlisting}

Then link your compiler binary:
\begin{lstlisting}[basicstyle=\ttfamily\small]
g++ my_compiler.cpp -lprecision_controller_ref -lmac_array_ref \
    -L /path/to/sw/reference_model -o my_compiler
\end{lstlisting}

\subsection{Shared Library / FFI}
\begin{lstlisting}[basicstyle=\ttfamily\small]
make shared    # produces *.so for FFI from any language
\end{lstlisting}

Python's \texttt{ctypes} or Rust's \texttt{libloading} can dlopen
the \texttt{.so} and call any \texttt{lhsi\_*} function declared in
\texttt{precision\_controller\_ref.hpp} or \texttt{mac\_array\_ref.hpp}.

\subsection{Direct Source Inclusion}
\begin{lstlisting}[basicstyle=\ttfamily\small]
g++ my_compiler.cpp \
    sw/reference_model/precision_controller_ref.cpp \
    sw/reference_model/mac_array_ref.cpp \
    -I sw/reference_model -o my_compiler
\end{lstlisting}

The reference \texttt{.cpp} files are small ($\sim$300 lines combined)
so direct inclusion is also fine.

%==============================================================================
\section{End-to-End Example}
\label{sec:e2e}

The cleanest pointer for a compiler engineer: run and read
\path{sw/reference_model/integration_example.py}. It composes both
blocks the way a compiler backend should --- process a $QK^{\top}$
tile, push to the precision controller, dispatch to either the INT8
or FP16 MAC path, accumulate.

\begin{lstlisting}[basicstyle=\ttfamily\small]
make integration
\end{lstlisting}

\noindent
Sample output (100 synthetic attention tiles, 21\% expected FP16):

\begin{lstlisting}[basicstyle=\ttfamily\small]
Routed:  INT8 = 82, FP16 = 18
Truth :  INT8 = 82, FP16 = 18

Per-tile cost:
  INT8:   1040 cycles, 131.1 nJ
  FP16:   4112 cycles, 655.4 nJ

Mixed-precision policy:       159k cycles, 22 uJ
All-FP16 baseline:            411k cycles, 66 uJ
Speedup:        2.58x
Energy saving:  65.6%
\end{lstlisting}

The routing matches ground truth exactly (82/18 vs 82/18); the cost
model shows the win from mixed-precision dispatch.

%==============================================================================
\section{Three Compiler-Use Scenarios}
\label{sec:scenarios}

\path{sw/reference_model/example_compiler_use.py} shows three levels
of sophistication for a compiler backend using just the precision
controller:

\begin{enumerate}
\item \textbf{Level A --- Naive per-tile dispatch.} For each tile, ask
      the gate and emit either an INT8 or FP16 kernel call. Simplest
      possible integration.
\item \textbf{Level B --- Batched index lists.} Collect decisions for
      a window of tiles, then batch all INT8 tiles to one kernel
      launch and all FP16 tiles to another. Amortizes per-tile
      overhead.
\item \textbf{Level C --- Offline calibration.} Run the gate over a
      calibration set offline, observe that some workloads are
      overwhelmingly one precision, and emit a workload-specific
      policy that skips the runtime gate.
\end{enumerate}

The end-to-end \texttt{integration\_example.py} (\S\ref{sec:e2e}) is
the next level up, composing both blocks.

%==============================================================================
\section{Verification Status}
\label{sec:verif}

\begin{table}[ht]
\centering
\small
\caption{Reference-model verification status.}
\label{tab:verif}
\begin{tabular}{@{}l l c@{}}
\toprule
Block                & Test                                                & Status \\
\midrule
Precision Controller & 143/143 RTL replay tiles, bit-exact (C++ and Py)    & pass \\
Precision Controller & Stateful streaming = batched = stateless            & pass \\
Precision Controller & extern "C" = C++ class                              & pass \\
Precision Controller & Canonical edges (spike $=$ 10 / 11)                 & pass \\
Precision Controller & Accumulator reset between tiles                     & pass \\
MAC Array            & INT8 matmul vs.\ naive int64 reference              & pass (exact) \\
MAC Array            & FP16 matmul vs.\ naive double reference             & pass (within $5 \cdot 10^{-3}$) \\
MAC Array            & Edge cases (zero, identity, signed extremes)        & pass \\
MAC Array            & extern "C" = C++ class                              & pass \\
MAC Array            & FP16 round-trip helpers                             & pass \\
MAC Array            & Cost-estimate sanity                                & pass \\
End-to-end           & Integration example routing matches ground truth    & 82/18 vs 82/18 \\
\bottomrule
\end{tabular}
\end{table}

%==============================================================================
\section{Where to Find Each File}
\label{sec:files}

\begin{lstlisting}[basicstyle=\ttfamily\small]
sw/reference_model/
  precision_controller_ref.hpp           -- block 1 header
  precision_controller_ref.cpp           -- block 1 implementation
  precision_controller_ref.py            -- block 1 Python ref
  test_precision_controller_ref.cpp      -- block 1 C++ tests
  test_precision_controller_ref.py       -- block 1 Python tests

  mac_array_ref.hpp                      -- block 2 header
  mac_array_ref.cpp                      -- block 2 implementation
  mac_array_ref.py                       -- block 2 Python ref
  test_mac_array_ref.cpp                 -- block 2 C++ tests
  test_mac_array_ref.py                  -- block 2 Python tests

  integration_example.py                 -- composes both blocks
  example_compiler_use.py                -- three sophistication levels

  Makefile                               -- build / test / shared / static
  README.md                              -- markdown form of this doc

docs/
  isa/precision_controller_isa.pdf       -- ISA spec for block 1
  mac_array_design.pdf                   -- design rationale for block 2
  sw_overview.pdf                        -- compiler-team entry point
  reference_model_api.pdf                -- THIS DOCUMENT
\end{lstlisting}

\vspace{0.5in}
\noindent\rule{\textwidth}{0.4pt}\\
\noindent\small\textit{This document is the canonical API reference
for the LonghornSilicon ACU reference models. Updated as each new
block lands in software.}

\end{document}
